\documentclass[10pt,a4paper]{amsart}
\usepackage[margin=19mm]{geometry}
\usepackage{amsmath,amssymb,mathtools}
\usepackage[scr=rsfs,cal=euler]{mathalfa}
\usepackage{xfrac}
\usepackage[unicode,hidelinks,bookmarks=false]{hyperref}
\newcommand{\ie}{i.e.\ }
\newcommand{\cf}{cf.\ }
\newcommand{\resp}{respectively\ }
\newcommand{\Subsection}{\subsection}
\providecommand{\nobreakdash}{\nobreak\hbox{-}}
\setlength{\emergencystretch}{2em}
\multlinegap0pt
\usepackage{bm}
\usepackage{bbm}
\usepackage{dsfont}
\usepackage{xspace}
\usepackage{cancel}
\usepackage{mathtools}
\usepackage{amsthm}
\makeatletter
\@ifundefined{lemma}{\newtheorem{lemma}{Lemma}}{}
\@ifundefined{proposition}{\newtheorem{proposition}{Proposition}}{}
\makeatother
\newcommand{\e}{\varepsilon}
\title[Time-asymptotic stability of viscous shocks for the outflow ]{Time-asymptotic stability of viscous shocks for the outflow problem of one-dimensional compressible fluids of Korteweg type}
\author{Sungho Han, Jeongho Kim, HyeonSeop Oh}
\date{Source: arXiv:2606.14247v1 (CC BY 4.0)}
\begin{document}
\maketitle

\noindent\emph{Source and licence.} Time-asymptotic stability of viscous shocks for the outflow problem of one-dimensional compressible fluids of Korteweg type. Version arXiv:2606.14247v1 (CC BY 4.0), distributed under the Creative Commons Attribution 4.0 International licence, as recorded in the arXiv licence metadata for this exact version and shown on the abstract page https://arxiv.org/abs/2606.14247v1. Full rights notice: licenses/arxiv-2606.14247-CC-BY-4.0.txt.\par\smallskip

\noindent\textit{Editorial scope.} Counted unit: the Lemma whose source label is \texttt{lem:higher2} in the article (source0.tex lines 1179--1188, source label \texttt{lem:higher2}), delivered together with the article's own complete original proof of it (source0.tex lines 1191--1292, one \texttt{proof} environment of 102 lines). No mathematics is written, repaired or re-derived: statement and proof are the article's own text line for line. Required context retained uncounted: eq:NSK (lines 109--114); cond:U+ (lines 160--162); prop:apriori (lines 363--409); p-assum (lines 375--377); small-X (lines 693--697); eq:perturbed-eq (lines 1010--1018); eq:GH-est (lines 1056--1058); est-quadratic (lines 1060--1066). Every definition, display and label of those lines is retained unchanged. Omitted from this package and not redistributed: 1--108, 115--159, 163--362, 378--692, 698--1009, 1019--1055, 1059--1059, 1067--1178, 1189--1190, 1293--1628. Nothing inside a retained excerpt is omitted or altered: every retained body is delivered exactly as the article's lines are written, so no omission receipt of any kind accompanies this package. Citations: the retained excerpts carry no citation command at all, so no bibliography entry of the article is transcribed and no reference is redistributed. Reference adaptation: none was needed and none was made; every reference inside the delivered unit resolves inside it, so no unresolved reference or citation is produced. Printed slips kept literal: no printed word, symbol, hypothesis or number of the counted statement or of its proof was altered. Layout and class: the article's own class is \texttt{amsart} with packages including geometry, amssymb, mathrsfs, amscd, bbm, extpfeil, graphicx, colortbl, float, caption, subcaption, comment, enumerate, url; the wrapper is the lane's pinned \texttt{amsart} editorial wrapper at 10pt on A4, which restates the theorem-like environments the retained text needs and adds the article's own macro definitions that the retained text needs (\texttt{\textbackslash e}), with extra wrapper packages (bm, bbm, dsfont, xspace, cancel, mathtools). The delivered numbering is the unit's own. Assets: no retained span contains a figure environment, a graphics-inclusion command, a PDF-page inclusion, a table or a TikZ picture; no image file is copied into this package, the package declares no asset dependency and no image file is redistributed with it. Licence: version arXiv:2606.14247v1 of this article is distributed under the Creative Commons Attribution 4.0 International licence, as recorded in the arXiv licence metadata for this exact version and shown on the abstract page https://arxiv.org/abs/2606.14247v1. A full rights notice is delivered at licenses/arxiv-2606.14247-CC-BY-4.0.txt. Characters: every retained excerpt is pure ASCII, so the delivered engine, XeLaTeX with its default Latin Modern text font, reports no missing character.
\begin{align} \label{eq:NSK}
\begin{aligned}
    &\rho_t +(\rho u)_x=0, \quad t>0, \quad x>0, \\
    &(\rho u)_t +(\rho u^2+p(\rho))_x = \mu u_{xx} + \kappa \rho\rho_{xxx},
\end{aligned}
\end{align}

\begin{equation}\label{cond:U+}
(\rho_+, u_+) \in \Omega_{sub}^- \cup  \Gamma_{trans}^-.
\end{equation}

\begin{proposition}\label{prop:apriori}

    For a given right-end state $(\rho_+, u_+) \in \mathbb{R}_+ \times \mathbb{R}$ satisfying $u_+<0$ and \eqref{cond:U+}, there exist positive constants $C_0, c_0, \delta_0$, and $\e_0$ such that the following holds.
    
    Let $u_- \in (u_+, 0)$ be the outflow velocity with shock strength $\delta := |u_- - u_+| < \delta_0$. Let $(\tilde{\rho}, \tilde{u})$ be the viscous-dispersive 2-shock profile connecting $(\rho_-, u_-)$ and $(\rho_+, u_+)$ with $\tilde{\rho}(0) = (\rho_- + \rho_+)/2$. Then, there exists a sufficiently large $\beta>0$, depending only on $\delta$, such that if $(\rho, u)$ is a strong solution to \eqref{eq:NSK} on $[0,T]$ satisfying
    \begin{equation*}
    \begin{aligned}
    &\rho - \tilde{\rho}^{X, \beta} \in C([0,T];H^2(\mathbb{R}_+)) \cap L^2(0,T;H^3(\mathbb{R}_+)),\\
    &u - \tilde{u}^{X, \beta} \in C([0,T];H^1(\mathbb{R}_+)) \cap L^2(0,T;H^2(\mathbb{R}_+)),
    \end{aligned} 
    \end{equation*}
    and
    \begin{equation}\label{p-assum}
    \sup_{t \in [0,T]} \Big( \| \rho - \tilde{\rho}^{X, \beta} \|_{H^2(\mathbb{R}_+)} + \| u - \tilde{u}^{X, \beta} \|_{H^1(\mathbb{R}_+)} \Big) \leq \e_0,
    \end{equation}
    then we have the following estimate for all $t \in [0, T]$:
    \begin{equation*}%\label{est:pri}
    \begin{aligned}
    & \|(\rho - \tilde{\rho}^{X, \beta})(t,\cdot
    )\|_{H^2(\mathbb{R}_+)}^2 + \| (u - \tilde{u}^{X, \beta})(t,\cdot) \|_{H^1(\mathbb{R}_+)}^2 \\
    &\qquad + \int_0^t \Big( \delta |\dot{X}|^2 + G_1 + G_3 + G^S + D_{u_1} + D_{u_2} + G_{w_1} + D_{w_1} + G_{w_2} + D_{w_2} \Big)\,ds\\
    &\quad \leq C_0 \Big( \| \rho_0 - \tilde{\rho}^{0,\beta} \|_{H^2(\mathbb{R}_+)}^2 + \| u_0 - \tilde{u}^{0,\beta} \|_{H^1(\mathbb{R}_+)}^2 \Big) + C_0 e^{-c_0 \delta \beta}.
    \end{aligned} 
    \end{equation*}
    Here, the constants $C_0$ and $c_0$ are independent of $T, \beta, \delta_0,$ and $\varepsilon_0$, and
    \begin{equation}\label{terms}
    \begin{split}
        &\begin{aligned}
         G_1 = \int_{\mathbb{R}_+} a_x \frac{\sigma-\tilde{u}}{2}\gamma\tilde{\rho}^{\gamma-2}\left((\rho-\tilde{\rho})-\frac{\tilde{\rho}}{\sigma-\tilde{u}}(u-\tilde{u})\right)^2\,dx,   
        \end{aligned}\\
        &\begin{aligned}
        &G_3 = \int_{\mathbb{R}_+} a_x \frac{\sigma-\tilde{u}}{2}\tilde{\rho}|w-\tilde{w}|^2\,dx, &&
        G^S = \int_{\mathbb{R}_+} |\tilde{u}_x| |u- \tilde{u}|^2\,dx, \\
        &D_{u_1} = \mu \int_{\mathbb{R}_+} a |(u - \tilde{u})_x|^2\,dx, &&
        D_{u_2} = \mu \int_{\mathbb{R}_+} \frac{1}{\rho} |(u - \tilde{u})_{xx}|^2\,dx, \\
        &G_{w_1} = \frac{1}{\sqrt{\kappa}} \int_{\mathbb{R}_+} \frac{p'(\rho)}{\rho^{1/2}} |w - \tilde{w}|^2\,dx, &&
        D_{w_1} = \sqrt{\kappa} \int_{\mathbb{R}_+} \rho^{1/2} |(w - \tilde{w})_x|^2\,dx, \\
        &G_{w_2} = \frac{1}{\sqrt{\kappa}} \int_{\mathbb{R}_+} \frac{p'(\rho)}{\rho^{1/2}} |(w - \tilde{w})_x|^2\,dx, &&
        D_{w_2} = \sqrt{\kappa} \int_{\mathbb{R}_+} \rho^{1/2} |(w - \tilde{w})_{xx}|^2\,dx.
        \end{aligned}
    \end{split}
    \end{equation}
    Furthermore, the derivative of the shift satisfies 
    \[
    |\dot{X}(t)| \leq C_0 \big( \| (\rho - \tilde{\rho}^{X, \beta})(t,\cdot) \|_{L^\infty(\mathbb{R}_+)} + \| (u - \tilde{u}^{X, \beta})(t,\cdot) \|_{L^\infty(\mathbb{R}_+)} \big), \quad \forall t \in [0,T].
    \]
\end{proposition}

\begin{equation}\label{small-X}
-\sigma t - X(t) - \beta \le -\frac{\sigma}{2}t - \beta < 0,
\qquad
\forall  t\le T.
\end{equation}

\begin{equation}\label{eq:perturbed-eq}
\left\{
\begin{aligned}
&\phi_t +u \phi_x +\rho \psi_x = F+\dot{X}(t) \tilde{\rho}_x,\\
&\rho\bigl(\psi_t+u \psi_x\bigr)+\frac{p'(\rho)\rho^{1/2}}{\sqrt{\kappa}}\omega =\mu \psi_{xx}+\sqrt{\kappa} (\rho^{3/2}\omega_x)_x +G+\dot{X}(t) \rho \tilde{u}_x,\\
&\rho(\omega_t +u \omega_x) = -\sqrt{\kappa} (\rho^{3/2}\psi_x)_x+ H+ \dot{X}(t) \rho \tilde{w}_x,
\end{aligned}
\right.
\end{equation}

\begin{equation} \label{eq:GH-est}
|(G, H)|\le C|\tilde{u}_x||(\phi,\psi,\phi_x)| \le C|\tilde{u}_x||(\phi,\psi,\omega)|.
\end{equation}

\begin{equation} \label{est-quadratic}
\begin{aligned}
\int_{\mathbb{R}_+} |\tilde{u}_x| |(\phi,\psi,\omega)|^2 \, dx 
&\le \int_{\mathbb{R}_+} \left[ \left( \phi-\frac{\tilde{\rho}}{\sigma-\tilde{u}}\psi\right)^2 + C |\psi|^2 + C|\omega|^2 \right] |\tilde{u}_x| \, dx \\
&\le C \sqrt{\delta} (G_1 + G_3) + G^S.
\end{aligned}
\end{equation}

\begin{lemma}\label{lem:higher2}
	Under the assumptions of Proposition \ref{prop:apriori}, for any $\tau \in (0,1)$, there exist positive constants $C, c$, and $C_\tau$ depending on $\tau$ such that, for all $t \in [0,T]$,
	\begin{equation} \label{est-higher-2}
		\begin{aligned}
			&\|(\psi_x, \omega_x)(t,\cdot)\|_{L^2}^2 + \int_0^t D_{u_2} \,ds \\
			&\le  \|(\psi_{0x}, \omega_{0x})\|_{L^2}^2 + \tau \int_0^t D_{w_2} \,ds + C\delta \int_0^t |\dot{X}(s)|^2 \,ds \\
			&\quad + C_\tau \int_0^t \left(G_1 + G_3 + G^S+ D_{u_1} + G_{w_1} + D_{w_1}   \right) ds + C e^{-c\delta \beta}.
		\end{aligned}
	\end{equation}
\end{lemma}

\begin{proof} We multiply $\eqref{eq:perturbed-eq}_2$ by $-\frac{\psi_{xx}}{\rho}$ and $\eqref{eq:perturbed-eq}_3$ by $-\frac{\omega_{xx}}{\rho}$, and then integrate the sum of the resulting two equations over $\mathbb{R}_+$ to get
\begin{equation}\label{eq:higher-order-1}
\begin{aligned}
    &-\int_{\mathbb{R}_+} (\psi_t \psi_{xx} + \omega_t \omega_{xx}) \,dx - \int_{\mathbb{R}_+} u (\psi_x \psi_{xx} + \omega_x \omega_{xx}) \,dx - \frac{1}{\sqrt{\kappa}}\int_{\mathbb{R}_+} \frac{p'(\rho)}{\rho^{1/2}} \omega \psi_{xx} \,dx \\
    &\quad + \mu \int_{\mathbb{R}_+} \frac{|\psi_{xx}|^2}{\rho}  \,dx + \sqrt{\kappa} \int_{\mathbb{R}_+} \frac{\psi_{xx}}{\rho} (\rho^{3/2} \omega_x)_x \,dx - \sqrt{\kappa} \int_{\mathbb{R}_+} \frac{\omega_{xx}}{\rho} (\rho^{3/2} \psi_x)_x \,dx \\
    &=- \int_{\mathbb{R}_+} \frac{G\psi_{xx} + H\omega_{xx}}{\rho} \,dx  - \dot{X}(t) \int_{\mathbb{R}_+} (\tilde{u}_x \psi_{xx} + \tilde{w}_x \omega_{xx}) \,dx.
\end{aligned}
\end{equation}

Applying integration by parts to the time derivative terms and directly expanding the terms including $\sqrt{\kappa}$ to cancel out the highest-order derivatives, \eqref{eq:higher-order-1} can be simplified to
\begin{equation*}%\label{eq:higher-order-1-1}
\begin{aligned}
    &\frac{1}{2} \frac{d}{dt} \|(\psi_x, \omega_x)\|_{L^2}^2 + \mu \int_{\mathbb{R}_+} \frac{|\psi_{xx}|^2}{\rho}  \,dx \\
    &= -(\psi_t \psi_x + \omega_t \omega_x)|_{x=0}   + \int_{\mathbb{R}_+} \left[u (\psi_x \psi_{xx} + \omega_x \omega_{xx}) +\frac{1}{\sqrt{\kappa}}\frac{p'(\rho)}{\rho^{1/2}} \omega \psi_{xx} \right] dx \\
    &\quad - \frac{3}{2}\sqrt{\kappa} \int_{\mathbb{R}_+} \frac{\rho_x}{\rho^{1/2}} (\psi_{xx} \omega_x - \omega_{xx} \psi_x) \,dx - \int_{\mathbb{R}_+} \frac{G\psi_{xx} + H\omega_{xx}}{\rho} \,dx \\
    &\quad - \dot{X}(t) \int_{\mathbb{R}_+} (\tilde{u}_x \psi_{xx} + \tilde{w}_x \omega_{xx}) \,dx \\
    &=: K_1 + K_2 + K_3 + K_4 + K_5.
\end{aligned}
\end{equation*}

For the boundary term $K_1$, the boundary conditions $u(t,0) = u_- < 0$ and $w(t,0) = 0$ and \eqref{small-X} imply
\begin{equation} \label{eq:psi_t_omega_t}
    \begin{aligned}
        |\psi_t(t,0)| &= |\tilde{u}_t(t,0)| \le C | \tilde{u}'(-\sigma t -X(t) -\beta)||\sigma + \dot{X}(t)|\le C \delta^2 e^{-c\delta ( \sigma t +\beta)},\\
        |\omega_t(t,0)| &= |\tilde{w}_t(t,0)| \le C | \tilde{w}'(-\sigma t -X(t) -\beta)||\sigma + \dot{X}(t)|\le C \delta^3 e^{-c\delta ( \sigma t +\beta)},\\
        |\omega(t,0)| &= |\tilde{w}(-\sigma t-X(t)-\beta)| \le C|\tilde{\rho}'(-\sigma t -X(t) -\beta)|\le C \delta^2 e^{-c\delta ( \sigma t +\beta)}.
    \end{aligned}
\end{equation}
Utilizing the decay estimates \eqref{eq:psi_t_omega_t} along with the Sobolev inequality, we obtain
\begin{equation*}
\begin{aligned}
|K_1| &\le \left( \| \psi_x\|_{L^\infty} |\psi_t(t,0)|  + \| \omega_x\|_{L^\infty} | \omega_t(t,0)|\right) \\
&\le C \|\psi_x\|_{L^\infty}^2|\psi_t(t,0)|+|\psi_t(t,0)| +\|\omega_x\|_{L^\infty}^2|\omega_t(t,0)|+ |\omega_t(t,0)|\\
&\le C  \delta^2\left( \|\psi_x\|_{L^2}^2+ \|\omega_x\|_{L^2}^2 +\|\psi_{xx}\|_{L^2}^2+\|\omega_{xx}\|_{L^2}^2 \right) + C\delta^2 e^{-c\delta(\sigma t+ \beta)} \\
&\le C \delta^2 (D_{u_1}+D_{w_1}+D_{u_2}+D_{w_2})+ C\delta^2 e^{-c\delta(\sigma t+ \beta)}.
\end{aligned}
\end{equation*}

%For $J_2$, we apply Young's inequality to obtain
%\begin{equation*}
%\begin{aligned}
%|J_2| &\le C \int_{\mathbb{R}_+} |\psi_x | | \psi_{xx} | \, dx + C \int_{\mathbb{R}_+} | \omega_x | | \omega_{xx} | \, dx + C \int_{\mathbb{R}_+} | \omega | | \psi_{xx} | \, dx \\
%&\le \frac{1}{8} D_{u_2} + \frac{1}{8} D_{w_2} + C (G_{w_1} + D_{u_1} + D_{w_1}).
%\end{aligned}
%\end{equation*}


For $K_2$, applying Young's inequality with a small parameter $\tau \in (0,1)$, we obtain
\begin{equation*}
\begin{aligned}
|K_2| &\le C \int_{\mathbb{R}_+} |\psi_x | | \psi_{xx} | \, dx + C \int_{\mathbb{R}_+} | \omega_x | | \omega_{xx} | \, dx + C \int_{\mathbb{R}_+} | \omega | | \psi_{xx} | \, dx \\
&\le \frac{1}{16} D_{u_2} + \frac{\tau}{4} D_{w_2} + C_\tau (G_{w_1} + D_{u_1} + D_{w_1}),
\end{aligned}
\end{equation*}
where the constant $C_\tau>0$ depends on $\tau$ but is independent of $T,\beta, \delta$ and $\varepsilon$.


For $K_3$, we recall \eqref{p-assum} implies $\|\rho_x\|_{L^\infty} \le C(\delta + \varepsilon)$, which yields
\begin{equation*}
\begin{aligned}
|K_3| &\le C \|\rho_x\|_{L^\infty} \int_{\mathbb{R}_+} \left( |\psi_{xx}||\omega_x| + |\omega_{xx}||\psi_x| \right) dx\le C (\delta + \varepsilon) \left( D_{u_1}+ D_{u_2} + D_{w_1} + D_{w_2} \right).
\end{aligned}
\end{equation*}

We use \eqref{eq:GH-est}, \eqref{est-quadratic} and Young's inequality to estimate $K_4$ as
\begin{equation*}
\begin{aligned}
|K_4| &\le \int_{\mathbb{R}_+}  |\tilde{u}_x| |(\phi,\psi,\omega)| |(\psi_{xx},\omega_{xx})| \, dx \\
&\le C\|\tilde{u}_x\|_{L^\infty}^{1/2} ( D_{u_2} +  D_{w_2} )+ C \|\tilde{u}_x\|_{L^\infty}^{1/2} \int_{\mathbb{R}_+} |\tilde{u}_x| |(\phi,\psi,\omega)|^2 \, dx \\
&\le C \delta \bigl(D_{u_2}+D_{w_2} +G_1+G_3 + G^S \bigr).
\end{aligned}
\end{equation*}

For $K_5$, we similarly obtain
\begin{equation*}
\begin{aligned}
|K_5| \le \delta |\dot{X}(t)|^2 + C \delta (D_{u_2} + D_{w_2} ).
\end{aligned}
\end{equation*}

%Combining the estimates from $J_1$ to $J_5$ and integrating over $[0,t]$, we use the smallness of $\delta$ and $\varepsilon$ to deduce
%\begin{equation} \label{est-higher-2}
%\begin{aligned}
%&\|(\psi_x, \omega_x)(t)\|_{L^2}^2 + \int_0^t D_{u_2} \,ds \\
%&\le  \|(\psi_x, \omega_x)(0)\|_{L^2}^2 + \frac{1}{2} \int_0^t D_{w_2} \,ds + \delta \int_0^t |\dot{X}(s)|^2 \,ds \\
%&\quad + C \int_0^t \left(G_1 + G_3 + G^S+ D_{u_1} + G_{w_1} + D_{w_1}   \right) ds + C \delta e^{-c\delta \beta},
%\end{aligned}
%\end{equation}
%which is the desired estimate.


Combining the estimates from $K_1$ to $K_5$ and integrating over $[0,t]$, we use the smallness of $\delta$ and $\varepsilon$ to deduce
\begin{equation*}
\begin{aligned}
&\|(\psi_x, \omega_x)(t,\cdot)\|_{L^2}^2 + \int_0^t D_{u_2} \,ds \\
&\le  \|(\psi_{0x}, \omega_{0x})\|_{L^2}^2 + \tau \int_0^t D_{w_2} \,ds + C\delta \int_0^t |\dot{X}(s)|^2 \,ds \\
&\quad + C_\tau \int_0^t \left(G_1 + G_3 + G^S+ D_{u_1} + G_{w_1} + D_{w_1}   \right) ds + C \delta e^{-c\delta \beta},
\end{aligned}
\end{equation*}
which is the desired estimate.

\end{proof}

\end{document}
